\section{Visual Prompting：概念与动机}

\subsection{什么是 Visual Prompting}

Visual Prompting（视觉提示）是一种在多模态视觉-语言模型（VLM）交互中，通过在原始图片上叠加视觉标记（如编号、箭头、掩码区域、边界框等），来辅助模型更好地进行视觉推理与定位的技巧。

\begin{importantbox}{Visual Prompting 的核心思想}
Visual Prompting 的本质是：在送给 VLM 之前，先用 CV 模块对原始图片做预处理（分割、检测、标注编号），再将增强后的图片送给 VLM 做推理。这样做将"直接生成坐标"的回归问题，转化为"选择编号"的分类问题，大幅降低了模型推理难度。
\end{importantbox}

以 Set-of-Mark（SoM）为例：给定一张原始图片和一个文本问题，SoM 首先用分割模型对图片做全图实例分割，然后为每个分割区域在其中心位置打上数字编号标签，最终将这张带编号的新图片连同文本问题一起交给 VLM 进行问答。

\begin{knowledgebox}{为什么 Visual Prompting 有效}
GPT-4V 等 VLM 在预训练阶段包含了大量带有标注的图表（labelled figures）、教科书插图（charts, textbooks）等图文对数据。这些训练数据中的图片本身就包含箭头、编号、分割区域等视觉标记。因此，当我们人为地在图片上添加类似标记时，能很自然地激活模型内部对这类标注的注意力机制，从而提升推理准确率。类似于数学几何题中的标注——编号、箭头、方向标记，都能帮助模型聚焦关键信息。
\end{knowledgebox}

\subsection{回归问题变分类问题}

Visual Prompting 最关键的优势在于降低推理难度：

\begin{itemize}
    \item \textbf{常规方式}：VLM 需要直接输出目标物体的精确坐标 $(x, y)$ 或边界框 $(x_1, y_1, x_2, y_2)$，这是一个回归问题
    \item \textbf{Visual Prompting 方式}：图片已经被分割并打上编号，VLM 只需选择正确的编号即可，这变成了一个分类/选择问题
\end{itemize}

\begin{warningbox}{Visual Prompting 并非万能}
随着现代 VLM（如 GPT-4o、Gemini 2.5 等）的视觉能力越来越强，部分场景下模型已经能直接准确输出坐标。但在复杂场景（多物体、遮挡、方位描述等）中，Visual Prompting 仍然是非常实用的辅助技巧。它通过帮助模型聚焦注意力并稳定推理过程，显著提升定位准确率。
\end{warningbox}

\subsection{本章小结}

Visual Prompting 是一种在原始图片上叠加 CV 预处理结果（分割掩码、数字编号、边界框等）的技巧，能将物体定位的回归问题转化为选择问题，降低 VLM 推理难度，提升准确率。其有效性源于 VLM 训练数据中天然包含大量带标注图片。


